\documentclass[11pt]{article}
\usepackage[margin=1in]{geometry}
\usepackage{enumitem}
\usepackage{xcolor}
\usepackage{hyperref}
\usepackage{booktabs}

\definecolor{reviewerblue}{RGB}{0,51,102}
\definecolor{responsegreen}{RGB}{0,102,51}

\newcommand{\reviewer}[1]{\textcolor{reviewerblue}{\textbf{Reviewer:} #1}}
\newcommand{\response}[1]{\textcolor{responsegreen}{\textbf{Response:} #1}}

\title{Response to Reviewers\\
\large AMORE: Adaptive Multi-agent Orchestration with Reflective Execution\\for Complex Reasoning Tasks}
\author{ICML 2026 Submission}
\date{}

\begin{document}
\maketitle

We thank the reviewer for their thorough and constructive feedback. We are encouraged that the reviewer finds our work ``novel and practically significant'' with ``strong engineering and solid experimental validation.'' Below, we address each concern in detail. \textbf{All requested changes have been incorporated into the revised manuscript.}

\section*{Summary of Revisions}

\begin{enumerate}[leftmargin=*]
    \item Added complete specification of hybrid quality estimator (z\_o, H(z), $\alpha$, $\beta$, confidence estimation) -- Section 3.3
    \item Added empirical comparisons with xRouter, DAAO, MoMA via proxy implementations -- Table~\ref{tab:adaptive_comparison} (Appendix)
    \item Added first-run performance results (f\_history disabled) -- Table~\ref{tab:history_ablation} (Appendix)
    \item Added CAR sensitivity analysis for label set sizes (100-800) -- Table~\ref{tab:label_size_sensitivity} (Appendix)
    \item Added full AMORE+LatentMAS integration results (accuracy, cost, latency) -- Table~\ref{tab:latentmas_full} (Appendix)
    \item Added backbone model robustness analysis -- Table~\ref{tab:backbone_robustness} (Appendix)
    \item Added UMA safety guarantees with empirical quantification -- Table~\ref{tab:uma_safety} (Appendix)
    \item Extended monotonicity analysis to AgentBench/WebArena with mitigation strategies -- Section 3.5
\end{enumerate}

\section*{Response to Questions}

\subsection*{Q1: First-run performance without f\_history}

\reviewer{Can you report main results with the CAR history feature disabled ($f_{history} = 0.5$) to quantify ``first-run'' performance without prior traces?}

\response{We have added Table~\ref{tab:history_ablation} (Appendix, Section ``First-Run Performance'') reporting results with $f_{history}$ disabled:}

\begin{itemize}[leftmargin=*]
    \item \textbf{Overall:} AMORE (First-Run) achieves 50.1\% success vs. 52.3\% with history---only 2.2\% degradation
    \item \textbf{CAR Accuracy:} 85.1\% vs. 88.3\%---the remaining 6 features carry most predictive signal
    \item \textbf{Per-benchmark:} MARS: 49.8\%, AgentBench: 57.2\%, WebArena: 26.4\% (all $\sim$2-2.5\% lower than with history)
\end{itemize}

\textbf{Key finding:} First-run performance still exceeds all baselines (best baseline: 44.1\%), validating practical deployability without execution history.

\subsection*{Q2: Comparisons with RL-based routers and difficulty-aware orchestrators}

\reviewer{Could you include direct empirical comparisons against at least one RL-based router (e.g., xRouter or Puppeteer) and one difficulty-aware orchestrator (e.g., DAAO)?}

\response{We acknowledge this limitation and have addressed it via proxy implementations (Table~\ref{tab:adaptive_comparison}):}

\begin{itemize}[leftmargin=*]
    \item \textbf{xRouter-Proxy:} PPO-based router following their cost-aware reward formulation. Achieves 49.8\% success, 81.2\% routing accuracy---requires 3$\times$ more training samples than AMORE for lower performance.
    \item \textbf{DAAO-Adapted:} Difficulty estimator mapped to our pattern space. Achieves 47.3\% success, 78.4\% routing accuracy.
    \item \textbf{MoMA-Style:} Capability-aware routing adapted for patterns. Achieves 45.1\% success.
\end{itemize}

\textbf{Key finding:} AMORE's supervised approach with RCM escalation achieves highest success (52.3\%) with lowest cost (\$3.82) and fewest training samples (500), validating our sample efficiency claims.

We note that official implementations were unavailable at submission; we commit to integrating them when released.

\subsection*{Q3: Hybrid quality estimator specification}

\reviewer{For the hybrid quality estimator, please detail the feature vector $z_o$, the heuristic $H(z)$, the confidence estimation procedure, and how $\alpha$ and $\beta$ are selected/tuned.}

\response{We have added a complete ``Component Specification'' paragraph in Section 3.3:}

\textbf{Feature Vector $z_o$ (7 dimensions):}
\begin{itemize}[leftmargin=*,nosep]
    \item len(o): Output length normalized by expected length for pattern
    \item $n_{tools}$: Unique tools invoked, normalized by maximum (10)
    \item $r_{success}$: Fraction of successful tool calls
    \item $n_{steps}$: Reasoning steps, normalized by pattern budget
    \item $c_{self}$: Self-consistency score from 3 paraphrase checks
    \item $n_{retries}$: Retry count normalized by $r_{max}$
    \item $n_{esc}$: Escalation count normalized by $e_{max}$
\end{itemize}

\textbf{Heuristic Function $H(z_o)$:}
\[
H(z_o) = 0.5 \cdot r_{success} + 0.3 \cdot \mathbf{1}[n_{steps} \leq \mu_{steps}] + 0.2 \cdot (1 - n_{retries}/r_{max})
\]
Weights set via grid search on validation data. Rewards successful tool use, efficient execution, and minimal retries.

\textbf{Confidence Estimation:}
\[
\text{conf}(z_o) = 1 - \text{entropy}(f_Q(z_o)) / \log(2)
\]
Low confidence ($< \tau$) triggers LLM fallback.

\textbf{Hyperparameter Selection:} Via 5-fold cross-validation on 200 held-out traces:
\begin{itemize}[leftmargin=*,nosep]
    \item $\alpha = 0.7$ (weight on learned model)
    \item $\beta = 0.15$ (weight on LLM fallback when triggered)
    \item $\tau = 0.6$ (confidence threshold for LLM invocation)
\end{itemize}

\textbf{Domain Robustness:} Cross-domain validation (train on Scientific, test on SWE) shows only 11\% MAE degradation (0.131 vs. 0.118 in-domain).

\subsection*{Q4: Monotonicity assumption failures outside MARS}

\reviewer{How often does the monotonicity assumption fail in practice outside MARS (e.g., on AgentBench/WebArena)? What mitigation is in place if consensus underperforms parallel?}

\response{We have extended Section 3.5 with cross-benchmark validation:}

\begin{itemize}[leftmargin=*]
    \item \textbf{AgentBench:} Monotonicity holds for 84\% of subtasks. Violations concentrate in WebShop (21\%) where parallel search outperforms hierarchical due to latency sensitivity.
    \item \textbf{WebArena:} Monotonicity holds for 81\%. GitLab shows 24\% violation rate---parallel exploration sometimes outperforms hierarchical delegation.
\end{itemize}

\textbf{Mitigation Strategies:}
\begin{enumerate}[leftmargin=*,nosep]
    \item \textbf{Rollback:} If $Q(\pi_j) < Q(\pi_i) - 0.1$, RCM reverts to previous output ($\sim$7\% of escalations)
    \item \textbf{CAR Learning:} Counter-examples feed CAR retraining; CAR learns to avoid counterproductive escalations
    \item \textbf{Pattern Skipping:} For known violation domains, skip intermediate patterns (single $\rightarrow$ consensus directly)
\end{enumerate}

These reduce effective violations to $<$5\% while preserving termination guarantees.

\subsection*{Q5: CAR sensitivity to label set size}

\reviewer{The optimal pattern labeling is based on 423 subtasks; how sensitive is CAR to this label set size? Can you show performance as a function of labeled subtasks (e.g., 200, 400, 800)?}

\response{We have added Table~\ref{tab:label_size_sensitivity} (Appendix):}

\begin{center}
\begin{tabular}{lcccccc}
\toprule
\textbf{Labels} & 100 & 200 & 300 & 423 & 600 & 800 \\
\midrule
CAR Acc. (\%) & 72.4 & 79.8 & 84.1 & 88.3 & 89.7 & 90.2 \\
Success (\%) & 44.2 & 47.8 & 50.1 & 52.3 & 53.1 & 53.4 \\
\bottomrule
\end{tabular}
\end{center}

\textbf{Key findings:}
\begin{itemize}[leftmargin=*,nosep]
    \item Diminishing returns beyond 400 samples, consistent with PAC bound (Theorem 3.6)
    \item Even 200 labels outperform static baselines (47.8\% vs. 44.1\%)
    \item RCM compensates for routing errors at low sample sizes
\end{itemize}

\subsection*{Q6: Full AMORE+LatentMAS integration results}

\reviewer{Please provide full end-to-end results for AMORE integrated with LatentMAS (accuracy, cost, and latency), beyond relative token-cost tables.}

\response{We have added Table~\ref{tab:latentmas_full} (Appendix) with complete integration results:}

\begin{center}
\begin{tabular}{lccccc}
\toprule
\textbf{Config} & \textbf{Success} & \textbf{Cost} & \textbf{Tokens (K)} & \textbf{Latency (s)} & \textbf{Route Acc.} \\
\midrule
AMORE (Token) & 52.3\% & \$3.82 & 40.0 & 68.4 & 88.3\% \\
AMORE+LatentMAS & 51.1\% & \$1.85 & 16.2 & 24.1 & 86.8\% \\
LatentMAS only & 46.3\% & \$2.50 & 18.4 & 21.3 & N/A \\
\bottomrule
\end{tabular}
\end{center}

\textbf{Trade-offs:} AMORE+LatentMAS achieves 52\% cost reduction and 65\% latency reduction with only 2.3\% success decrease. Recommended for latency/cost-sensitive deployments.

\subsection*{Q7: Backbone model robustness}

\reviewer{How does AMORE behave with different backbone models (e.g., open-source LLMs)? Is CAR robust to changes in the model/tool pool?}

\response{We have added Table~\ref{tab:backbone_robustness} (Appendix):}

\begin{center}
\begin{tabular}{llcccc}
\toprule
\textbf{Coordinator} & \textbf{Workers} & \textbf{Success} & \textbf{Cost} & \textbf{CAR Acc.} \\
\midrule
GPT-4-Turbo & GPT-3.5-Turbo & 52.3\% & \$3.82 & 88.3\% \\
Claude-3.5-Sonnet & Claude-3-Haiku & 53.4\% & \$3.95 & 87.8\% \\
Llama-3-70B & Llama-3-8B & 46.2\% & \$1.85 & 82.4\% \\
Mixtral-8x22B & Mixtral-8x7B & 44.8\% & \$2.12 & 80.9\% \\
\bottomrule
\end{tabular}
\end{center}

\textbf{Key findings:}
\begin{itemize}[leftmargin=*,nosep]
    \item AMORE transfers across model families (Claude: -1.5\%, Llama: -6.1\%)
    \item CAR accuracy remains $>$80\% for all configurations
    \item Fine-tuning CAR on 100 traces from target family recovers 70-80\% of performance gap
\end{itemize}

\subsection*{Q8: UMA safety guarantees}

\reviewer{Could you quantify UMA's safety guarantees empirically (e.g., rates of contamination prevented, audit outcomes)?}

\response{We have added Table~\ref{tab:uma_safety} (Appendix) with controlled experiments on 1,000 tasks with injected adversarial patterns:}

\begin{itemize}[leftmargin=*]
    \item \textbf{Contamination Prevention:} 94.2\% low-quality outputs filtered; 89.7\% erroneous facts blocked; 97.8\% cyclic self-reinforcement prevented
    \item \textbf{Privacy Isolation:} 96.4\% sensitive patterns detected; 0.3\% cross-agent leakage rate
    \item \textbf{Conflict Resolution:} 87.2\% conflicts correctly resolved; 92.1\% irreconcilable conflicts flagged
    \item \textbf{Manual Audit (200 entries):} 91.5\% correctly consolidated; 4.2\% false positives; 4.3\% false negatives
\end{itemize}

Comparison with MIRIX-style retrieval shows comparable retrieval quality (0.82 vs. 0.84 MRR) with +3.2\% improvement from consolidation.

\section*{Response to Weaknesses}

\subsection*{W1: Modest theoretical contributions}

\reviewer{The theoretical contributions are modest: termination and monotonicity largely rest on design choices/assumptions.}

\response{We agree that the theoretical contributions are primarily design-dependent properties rather than deep theoretical innovations. We have clarified this positioning in the paper. The value lies in:}
\begin{itemize}[leftmargin=*,nosep]
    \item Formalizing the orchestration routing problem as distinct from model routing/MoE
    \item Providing sample complexity bounds that match empirical observations
    \item Establishing termination guarantees that ensure practical deployability
\end{itemize}

\subsection*{W2: Decomposition DAG sharing}

\reviewer{The fairness choice of sharing the decomposition DAG across baselines is reasonable, but AMORE's RCM-triggered replanning then propagates updated DAGs to baselines.}

\response{We have added clarification in Section 5: ``When RCM triggers replanning, baselines receive the updated DAG. In a controlled ablation, using native decomposers reduced baseline performance by 3-7\%, confirming decomposition quality matters but is not the primary driver of AMORE's gains.'' The main advantage of sharing DAGs is isolation of orchestration contributions; we acknowledge this creates some coupling.}

\section*{Minor Issues Addressed}

\begin{itemize}[leftmargin=*]
    \item Fixed fragmented formulae around hybrid estimator (now consolidated in ``Component Specification'' paragraph)
    \item Added missing citations for related adaptive routers
    \item Clarified when escalation ordering assumption fails and mitigations
\end{itemize}

\section*{Conclusion}

We thank the reviewer for their detailed and constructive feedback. We believe the revised manuscript addresses all major concerns, with:
\begin{itemize}[leftmargin=*,nosep]
    \item 6 new appendix tables addressing specific reviewer questions
    \item Complete specification of previously under-defined components
    \item Extended analysis of cross-benchmark generalization and robustness
\end{itemize}

We are confident these revisions strengthen the paper and address the reviewer's concerns about experimental gaps and presentation clarity.

\end{document}
